% Colour marks: \mk{<colour id>}{<text>} highlights a statement of the model text in the colour of the
% passage that supports it in the cited paper (sources_marked/<module>/<paper>.pdf).
% The colour ids and RGB values are those of tools/mark_pdfs.py; the claim behind each id is listed
% per module in tools/marks/<module>.json.  Remove the % Colour marks: \mk{<colour id>}{<text>} highlights a statement of the model text in the colour of the
% passage that supports it in the cited paper (sources_marked/<module>/<paper>.pdf).
% The colour ids and RGB values are those of tools/mark_pdfs.py; the claim behind each id is listed
% per module in tools/marks/<module>.json.  Remove the % Colour marks: \mk{<colour id>}{<text>} highlights a statement of the model text in the colour of the
% passage that supports it in the cited paper (sources_marked/<module>/<paper>.pdf).
% The colour ids and RGB values are those of tools/mark_pdfs.py; the claim behind each id is listed
% per module in tools/marks/<module>.json.  Remove the % Colour marks: \mk{<colour id>}{<text>} highlights a statement of the model text in the colour of the
% passage that supports it in the cited paper (sources_marked/<module>/<paper>.pdf).
% The colour ids and RGB values are those of tools/mark_pdfs.py; the claim behind each id is listed
% per module in tools/marks/<module>.json.  Remove the \input{marks} line to typeset without marks.
\definecolor{mk1}{RGB}{255,235,59}   % yellow
\definecolor{mk2}{RGB}{139,220,110}  % green
\definecolor{mk3}{RGB}{100,215,240}  % cyan
\definecolor{mk4}{RGB}{255,150,200}  % pink
\definecolor{mk5}{RGB}{255,178,80}   % orange
\definecolor{mk6}{RGB}{200,160,255}  % violet
\definecolor{mk7}{RGB}{215,180,130}  % tan
\definecolor{mk8}{RGB}{140,170,255}  % blue
\definecolor{mk9}{RGB}{255,120,110}  % salmon
\definecolor{mk10}{RGB}{190,190,190} % grey
\newcommand{\mk}[2]{\begingroup\sethlcolor{mk#1}\hl{#2}\endgroup}
% \mkt: same for text inside tabularx cells (soul cannot be used there); the marked text does not break lines.
\newcommand{\mkt}[2]{{\setlength{\fboxsep}{0.5pt}\colorbox{mk#1}{#2}}}
 line to typeset without marks.
\definecolor{mk1}{RGB}{255,235,59}   % yellow
\definecolor{mk2}{RGB}{139,220,110}  % green
\definecolor{mk3}{RGB}{100,215,240}  % cyan
\definecolor{mk4}{RGB}{255,150,200}  % pink
\definecolor{mk5}{RGB}{255,178,80}   % orange
\definecolor{mk6}{RGB}{200,160,255}  % violet
\definecolor{mk7}{RGB}{215,180,130}  % tan
\definecolor{mk8}{RGB}{140,170,255}  % blue
\definecolor{mk9}{RGB}{255,120,110}  % salmon
\definecolor{mk10}{RGB}{190,190,190} % grey
\newcommand{\mk}[2]{\begingroup\sethlcolor{mk#1}\hl{#2}\endgroup}
% \mkt: same for text inside tabularx cells (soul cannot be used there); the marked text does not break lines.
\newcommand{\mkt}[2]{{\setlength{\fboxsep}{0.5pt}\colorbox{mk#1}{#2}}}
 line to typeset without marks.
\definecolor{mk1}{RGB}{255,235,59}   % yellow
\definecolor{mk2}{RGB}{139,220,110}  % green
\definecolor{mk3}{RGB}{100,215,240}  % cyan
\definecolor{mk4}{RGB}{255,150,200}  % pink
\definecolor{mk5}{RGB}{255,178,80}   % orange
\definecolor{mk6}{RGB}{200,160,255}  % violet
\definecolor{mk7}{RGB}{215,180,130}  % tan
\definecolor{mk8}{RGB}{140,170,255}  % blue
\definecolor{mk9}{RGB}{255,120,110}  % salmon
\definecolor{mk10}{RGB}{190,190,190} % grey
\newcommand{\mk}[2]{\begingroup\sethlcolor{mk#1}\hl{#2}\endgroup}
% \mkt: same for text inside tabularx cells (soul cannot be used there); the marked text does not break lines.
\newcommand{\mkt}[2]{{\setlength{\fboxsep}{0.5pt}\colorbox{mk#1}{#2}}}
 line to typeset without marks.
\definecolor{mk1}{RGB}{255,235,59}   % yellow
\definecolor{mk2}{RGB}{139,220,110}  % green
\definecolor{mk3}{RGB}{100,215,240}  % cyan
\definecolor{mk4}{RGB}{255,150,200}  % pink
\definecolor{mk5}{RGB}{255,178,80}   % orange
\definecolor{mk6}{RGB}{200,160,255}  % violet
\definecolor{mk7}{RGB}{215,180,130}  % tan
\definecolor{mk8}{RGB}{140,170,255}  % blue
\definecolor{mk9}{RGB}{255,120,110}  % salmon
\definecolor{mk10}{RGB}{190,190,190} % grey
\newcommand{\mk}[2]{\begingroup\sethlcolor{mk#1}\hl{#2}\endgroup}
% \mkt: same for text inside tabularx cells (soul cannot be used there); the marked text does not break lines.
\newcommand{\mkt}[2]{{\setlength{\fboxsep}{0.5pt}\colorbox{mk#1}{#2}}}
